% =====================================================================
%  Bölüm 5 — Doğal Çıkarım Ağaçları ve Lean Taktikleri
% =====================================================================
\chapter{Doğal Çıkarım Ağaçları ve Lean Taktikleri}\label{bol:dogal-cikarim}

\begin{bolumhedefleri}
\begin{itemize}
  \item Önermeler mantığının doğal çıkarım kurallarının tamamını tek bir
        tabloda görmeyi,
  \item \Lean{}'in hedef ekranının bir \emph{ardışık} ($\Gamma \vdash \varphi$)
        olduğunu ve her taktiğin bir ardışığı başkalarına dönüştürdüğünü,
  \item Geriye doğru (\leantac{apply}, \leantac{suffices}) ve ileriye doğru
        (\leantac{have}) akıl yürütmeyi ve ikisini birlikte kullanmayı,
  \item Bir ispata nereden başlanacağını söyleyen bir strateji haritasını,
  \item İspat normalleştirmeyi: dolambaçlı ispatların nasıl
        sadeleştirildiğini ve bunun neden ``hesaplama'' olduğunu,
  \item Taktik birleştiricilerini: \lc{<;>}, \leantac{all\_goals},
        \leantac{first}, \leantac{refine}.
\end{itemize}
\end{bolumhedefleri}

Bu bölüm Kısım~I'i kapatıyor. Yeni bir bağlaç tanıtmayacağız; bunun
yerine önceki dört bölümde öğrendiklerimizi bir araya getirip
\emph{düzenleyeceğiz}. Amaç, önünüze yeni bir önerme geldiğinde
``nereden başlamalıyım?'' sorusunu sistemli biçimde cevaplayabilmenizdir.

% ---------------------------------------------------------------------
\section{Kural kataloğu}\label{sec:kural-katalogu}

\Cref{fig:kural-katalogu}, Bölüm~2--4'te tanıştığımız bütün çıkarım
kurallarını bir arada gösteriyor. Tabloda iki şeye dikkat edin.
Birincisi, her bağlacın bir \textcolor{ikincil}{\bfseries giriş} ve bir
\textcolor{vurgu}{\bfseries eleme} sütunu vardır. İkincisi, yalnızca son
satırdaki kural (\emph{klasik saçmaya indirgeme}) yapıcı değildir.

\begin{figure}[p]
\centering
\small
\renewcommand{\arraystretch}{1}
\setlength{\tabcolsep}{4pt}
\begin{tabular}{>{\centering\arraybackslash}m{1.2cm}
                >{\centering\arraybackslash}m{5.3cm}
                >{\centering\arraybackslash}m{6.6cm}}
\rowcolor{tabloust}
\tablobaslik{Bağlaç} & \tablobaslik{Giriş (kanıtlamak)} & \tablobaslik{Eleme (kullanmak)} \\[2pt]
% →
$\to$ &
\vspace{2pt}
\begin{prooftree}
  \hypo{[P]^{n}} \ellipsis{}{Q} \infer1[$\to$G$^n$]{P \to Q}
\end{prooftree}\par\smallskip \leantac{intro} &
\vspace{2pt}
\begin{prooftree}
  \hypo{P \to Q} \hypo{P} \infer2[$\to$E]{Q}
\end{prooftree}\par\smallskip \leantac{apply}, \lc{h hp} \\[4pt]
\hline
% ∧
$\wedge$ &
\vspace{2pt}
\begin{prooftree}
  \hypo{P} \hypo{Q} \infer2[$\wedge$G]{P \wedge Q}
\end{prooftree}\par\smallskip \leantac{constructor}, \lc{⟨hp, hq⟩} &
\vspace{2pt}
\begin{prooftree}
  \hypo{P \wedge Q} \infer1[$\wedge$E$_1$]{P}
\end{prooftree}\quad
\begin{prooftree}
  \hypo{P \wedge Q} \infer1[$\wedge$E$_2$]{Q}
\end{prooftree}\par\smallskip \leantac{obtain}, \lc{h.1}, \lc{h.2} \\[4pt]
\hline
% ∨
$\vee$ &
\vspace{2pt}
\begin{prooftree}
  \hypo{P} \infer1[$\vee$G$_1$]{P \vee Q}
\end{prooftree}\quad
\begin{prooftree}
  \hypo{Q} \infer1[$\vee$G$_2$]{P \vee Q}
\end{prooftree}\par\smallskip \lc{Or.inl}, \lc{Or.inr} &
\vspace{2pt}
\begin{prooftree}
  \hypo{P \vee Q}
  \hypo{[P]^{m}} \ellipsis{}{R}
  \hypo{[Q]^{n}} \ellipsis{}{R}
  \infer3[$\vee$E$^{m,n}$]{R}
\end{prooftree}\par\smallskip \leantac{cases}, \leantac{rcases} \\[4pt]
\hline
% ¬
$\neg$ &
\vspace{2pt}
\begin{prooftree}
  \hypo{[P]^{n}} \ellipsis{}{\bot} \infer1[$\neg$G$^n$]{\neg P}
\end{prooftree}\par\smallskip \leantac{intro} &
\vspace{2pt}
\begin{prooftree}
  \hypo{\neg P} \hypo{P} \infer2[$\neg$E]{\bot}
\end{prooftree}\par\smallskip \lc{hnp hp} \\[4pt]
\hline
% ↔
$\leftrightarrow$ &
\vspace{2pt}
\begin{prooftree}
  \hypo{[P]^{m}} \ellipsis{}{Q}
  \hypo{[Q]^{n}} \ellipsis{}{P}
  \infer2[$\leftrightarrow$G$^{m,n}$]{P \leftrightarrow Q}
\end{prooftree}\par\smallskip \leantac{constructor} &
\vspace{2pt}
\begin{prooftree}
  \hypo{P \leftrightarrow Q} \hypo{P} \infer2[$\leftrightarrow$E$_1$]{Q}
\end{prooftree}\quad
\begin{prooftree}
  \hypo{P \leftrightarrow Q} \hypo{Q} \infer2[$\leftrightarrow$E$_2$]{P}
\end{prooftree}\par\smallskip \lc{h.mp}, \lc{h.mpr}, \leantac{rw} \\[4pt]
\hline
% ⊤ ⊥
$\top$, $\bot$ &
\vspace{2pt}
\begin{prooftree}
  \infer0[$\top$G]{\top}
\end{prooftree}\par\smallskip \leantac{trivial} &
\vspace{2pt}
\begin{prooftree}
  \hypo{\bot} \infer1[$\bot$E]{P}
\end{prooftree}\par\smallskip \leantac{exfalso}, \lc{False.elim} \\[4pt]
\hline
\rowcolor{teorembg}
\textcolor{teoremframe}{\bfseries klasik} &
\multicolumn{2}{c}{%
\vspace{2pt}
\begin{prooftree}
  \hypo{[\neg P]^{n}} \ellipsis{}{\bot} \infer1[RAA$^n$]{P}
\end{prooftree}\quad\raisebox{1.2em}{\lc{Classical.byContradiction}}} \\[4pt]
\end{tabular}

\vspace{6pt}
\begin{tikzpicture}
  \node[draw=black!20, rounded corners=3pt, fill=white, font=\footnotesize,
        text width=12.6cm, align=left, inner sep=6pt] {%
    \textbf{Varsayım kuralı.} Bir varsayım, kendi başına bir türetimdir:
    \quad$\begin{prooftree}\infer0{P}\end{prooftree}$ \;(varsayım $P$).
    \Lean{}'de \leantac{assumption} ya da \leantac{exact hp}.
    \quad \textbf{Kapatma.} $[P]^n$ biçimindeki varsayımlar, aynı
    numarayı taşıyan kuralda kapatılır.};
\end{tikzpicture}
\caption{Önermeler mantığı için doğal çıkarımın kural kataloğu ve her
kuralın \Lean{} karşılığı. Pembe satır dışındaki bütün kurallar
yapıcıdır; pembe satır klasik mantığa geçiştir.}
\label{fig:kural-katalogu}
\end{figure}

\begin{tanim}{Doğal çıkarım türetimi}{turetim}
Bir \textbf{türetim}, yaprakları varsayımlar, iç düğümleri
\cref{fig:kural-katalogu}'deki kuralların uygulamaları olan bir
ağaçtır. Kökteki önerme türetimin \textbf{sonucudur}. Bir türetimin
\textbf{açık varsayımları}, hiçbir kuralda kapatılmamış yapraklardır.
$\Gamma$ açık varsayımlar kümesi ve $\varphi$ sonuç ise,
\[
  \Gamma \vdash \varphi
\]
yazarız ve ``$\varphi$, $\Gamma$'dan türetilebilir'' diye okuruz.
Açık varsayımı olmayan bir türetimin sonucuna \textbf{teorem} denir:
$\vdash \varphi$.
\end{tanim}

% ---------------------------------------------------------------------
\section{Lean'in hedef ekranı bir ardışıktır}\label{sec:ardisik-gorunum}

Bölüm~2'den beri her \leantac{intro} adımında varsayımların Infoview
panelinde \lc{⊢} işaretinin \emph{üstüne} taşındığını gördük. Bu
gösterim bir rastlantı değildir: \Lean{}'in hedef ekranı, tam olarak
$\Gamma \vdash \varphi$ biçiminde bir \emph{ardışıktır} (İng.\
\emph{sequent}) (\cref{fig:ardisik}).

\begin{figure}[htbp]
\centering
\begin{tikzpicture}[
    kutu/.style={draw=hedefframe, fill=hedefbg, thick, rounded corners=3pt,
                 font=\leanfont\small, align=left, inner sep=7pt},
    ok/.style={-{Stealth[length=2.6mm]}, thick, #1},
  ]
  \node[kutu] (lean) at (0,0) {%
    P Q : \textcolor{leantype}{\bfseries Prop}\\
    \textcolor{leanhyp}{\bfseries hp} : P\\
    \textcolor{leanhyp}{\bfseries hq} : Q\\[3pt]
    \textcolor{leansym}{⊢} P \textcolor{leansym}{∧} Q};
  \node[font=\Large] (es) at (3.2,0) {$\Longleftrightarrow$};
  \node[font=\Large] (sq) at (6.6,0) {$\underbrace{P,\; Q}_{\Gamma} \;\vdash\; \underbrace{P \wedge Q}_{\varphi}$};
  \draw[decorate, decoration={brace, amplitude=5pt}, thick, leanhyp]
    ([xshift=4pt]lean.north east) -- ([xshift=4pt,yshift=-1.55cm]lean.north east)
    node[midway, right=7pt, font=\baslikfont\scriptsize, text=leanhyp]{bağlam $\Gamma$};
  \node[below=0.35cm of lean, font=\baslikfont\scriptsize, text=black!60] {Lean Infoview};
  \node[below=0.55cm of sq, font=\baslikfont\scriptsize, text=black!60] {ardışık (sequent)};
\end{tikzpicture}
\caption{\Lean{}'in hedef durumu ile mantıktaki ardışık gösterimi aynı
şeydir. Hipotez adları (\lc{hp}, \lc{hq}), varsayımların numaralarının
yerini tutar.}
\label{fig:ardisik}
\end{figure}

Bu bakış açısıyla doğal çıkarım kurallarını, varsayımları ağacın
yapraklarında taşımak yerine her düğümde açıkça yazarak ifade
edebiliriz. Örneğin varsayım kuralı ile implikasyonun kuralları şöyle
olur:

\begin{center}
\begin{tabular}{c@{\hspace{2.6em}}c@{\hspace{2.6em}}c}
\begin{prooftree}
  \infer0[Var]{\Gamma, P \vdash P}
\end{prooftree}
&
\begin{prooftree}
  \hypo{\Gamma, P \vdash Q}
  \infer1[$\to$G]{\Gamma \vdash P \to Q}
\end{prooftree}
&
\begin{prooftree}
  \hypo{\Gamma \vdash P \to Q}
  \hypo{\Gamma \vdash P}
  \infer2[$\to$E]{\Gamma \vdash Q}
\end{prooftree}
\\[10pt]
\small \leantac{assumption} & \small \leantac{intro} & \small \leantac{apply}
\end{tabular}
\end{center}

\noindent Bu biçimde okunduğunda her kural bir \emph{taktiğin tarifidir}:

\begin{onemlifikir}[Önemli Fikir: Taktik = kuralı aşağıdan yukarı uygulamak]
Bir çıkarım kuralı, \emph{yukarıdaki} ardışıklardan \emph{aşağıdakine}
geçişi gösterir. Bir taktik ise aynı kuralı \emph{tersten} uygular:
Aşağıdaki ardışığı (mevcut hedefi) alır ve onu yukarıdaki ardışıklara
(yeni hedeflere) dönüştürür. \leantac{intro}, $\Gamma \vdash P \to Q$
hedefini $\Gamma, P \vdash Q$ hedefine çevirir. \leantac{constructor},
$\Gamma \vdash P \wedge Q$ hedefini iki hedefe böler. Kâğıt üzerinde ağacı
yapraklardan köke doğru \emph{okuruz}; \Lean{}'de ise kökten yapraklara
doğru \emph{inşa ederiz}.
\end{onemlifikir}

\leandosya[Varsayım kuralı]{Bolum05}{varsayim}

\noindent Ardışık biçimindeki bu doğal çıkarım sistemi, \cref{bol:ardisik}'te
göreceğimiz Gentzen'in \emph{ardışık hesabından} farklıdır: Orada sağ
tarafta da birden çok önerme bulunabilir ve eleme kuralları yerine
``sol giriş'' kuralları kullanılır.

% ---------------------------------------------------------------------
\section{Kural--taktik eşleme tablosu}\label{sec:kural-taktik}

\Cref{tab:kural-taktik}, her kuralın taktik modunda (geriye doğru) ve
terim modunda (ileriye doğru) nasıl kullanıldığını özetliyor. Bu tablo,
kitabın geri kalanında sık sık başvuracağınız bir kaynak olacak.

\begin{table}[htbp]
\centering
\small
\renewcommand{\arraystretch}{1.45}
\rowcolors{2}{tablosatir}{white}
\begin{tabularx}{\linewidth}{>{\raggedright\arraybackslash}p{1.6cm}
                             >{\raggedright\arraybackslash}p{3.1cm}
                             >{\raggedright\arraybackslash}X
                             >{\centering\arraybackslash}p{1.0cm}}
\rowcolor{tabloust}
\tablobaslik{Kural} & \tablobaslik{Taktik (hedef üzerinde)} & \tablobaslik{Terim / hipotez üzerinde} & \tablobaslik{Bölüm} \\
Var                 & \leantac{assumption}, \leantac{exact h} & \lc{h} & \ref{bol:giris} \\
$\to$G              & \leantac{intro h}             & \lc{fun h => t} & \ref{bol:onermeler1} \\
$\to$E              & \leantac{apply f}             & \lc{f a} & \ref{bol:onermeler1} \\
$\wedge$G           & \leantac{constructor}         & \lc{⟨a, b⟩} & \ref{bol:onermeler1} \\
$\wedge$E           & ---                           & \lc{h.1}, \lc{h.2}, \leantac{obtain ⟨a, b⟩ := h} & \ref{bol:onermeler1} \\
$\vee$G             & \leantac{apply Or.inl}        & \lc{Or.inl a}, \lc{Or.inr b} & \ref{bol:onermeler1} \\
$\vee$E             & ---                           & \leantac{cases h}, \leantac{rcases h with a | b}, \lc{Or.elim} & \ref{bol:onermeler1} \\
$\neg$G             & \leantac{intro h}             & \lc{fun h => ...} & \ref{bol:degilleme} \\
$\neg$E             & \leantac{apply hn}            & \lc{hn h}, \lc{absurd h hn} & \ref{bol:degilleme} \\
$\leftrightarrow$G  & \leantac{constructor}         & \lc{⟨f, g⟩} & \ref{bol:degilleme} \\
$\leftrightarrow$E  & \leantac{rw [h]}              & \lc{h.mp a}, \lc{h.mpr b} & \ref{bol:degilleme} \\
$\top$G             & \leantac{trivial}             & \lc{True.intro} & \ref{bol:degilleme} \\
$\bot$E             & \leantac{exfalso}             & \lc{False.elim h}, \leantac{cases h} & \ref{bol:degilleme} \\
RAA                 & \leantac{apply Classical.}\newline\leantac{\quad byContradiction} & \lc{Classical.byContradiction f} & \ref{bol:yapici-klasik} \\
\end{tabularx}
\caption{Doğal çıkarım kuralları ile \Lean{} taktiklerinin ve terimlerinin
eşlemesi. Eleme kuralları çoğunlukla \emph{hipotezler} üzerinde, giriş
kuralları ise \emph{hedef} üzerinde çalışır.}
\label{tab:kural-taktik}
\end{table}

% ---------------------------------------------------------------------
\section{Geriye ve ileriye akıl yürütme}\label{sec:geri-ileri}

Bir ispatı iki yönden kurabiliriz (\cref{fig:geri-ileri}):

\begin{itemize}
  \item \textbf{Geriye doğru}: Hedeften başlar, ``bunu göstermek için ne
        yeter?'' diye sorarız. \leantac{apply}, \leantac{intro},
        \leantac{constructor}, \leantac{suffices} bu yönde çalışır.
  \item \textbf{İleriye doğru}: Hipotezlerden başlar, ``bunlardan ne
        çıkar?'' diye sorarız. \leantac{have} ve terim modundaki
        uygulamalar bu yönde çalışır.
\end{itemize}

\begin{figure}[htbp]
\centering
\begin{tikzpicture}[
    dugum/.style={draw=#1, fill=#1!8, thick, rounded corners=3pt,
                  minimum width=1.5cm, minimum height=0.8cm, font=\leanfont\small},
    ok/.style={-{Stealth[length=2.6mm]}, very thick, #1},
  ]
  % hipotezler
  \node[dugum=leanhyp] (p) at (0,0) {p : P};
  \node[dugum=leanframe] (q) at (3.6,0) {Q};
  \node[dugum=leanframe] (r) at (7.2,0) {R};
  \node[dugum=bittiframe] (g) at (10.6,0) {⊢ R};
  \draw[ok=leansym] (p) -- node[above, font=\leanfont\scriptsize\bfseries, text=leanhyp]{q := h1 p} (q);
  \draw[ok=leansym] (q) -- node[above, font=\leanfont\scriptsize\bfseries, text=leanhyp]{r := h2 q} (r);
  \draw[ok=leansym, dashed] (r) -- node[above, font=\leanfont\scriptsize\bfseries, text=leanhyp]{exact r} (g);
  \node[font=\baslikfont\small\bfseries, text=leansym, anchor=west] at (-0.8,0.95) {İleriye doğru (have)};
  % geri
  \node[dugum=hedefframe] (g1) at (10.6,-2.3) {⊢ R};
  \node[dugum=hedefframe] (g2) at (7.2,-2.3) {⊢ Q};
  \node[dugum=hedefframe] (g3) at (3.6,-2.3) {⊢ P};
  \node[dugum=bittiframe] (g4) at (0,-2.3) {p : P};
  \draw[ok=leantac] (g1) -- node[below, font=\leanfont\scriptsize\bfseries, text=leantac]{apply h2} (g2);
  \draw[ok=leantac] (g2) -- node[below, font=\leanfont\scriptsize\bfseries, text=leantac]{apply h1} (g3);
  \draw[ok=leantac, dashed] (g3) -- node[below, font=\leanfont\scriptsize\bfseries, text=leantac]{exact p} (g4);
  \node[font=\baslikfont\small\bfseries, text=leantac, anchor=east] at (11.4,-3.25) {Geriye doğru (apply)};
\end{tikzpicture}
\caption{Aynı ispat iki yönden: İleriye doğru ispat hipotezlerden
hedefe, geriye doğru ispat hedeften hipotezlere yürür. İkisi ortada
buluşur.}
\label{fig:geri-ileri}
\end{figure}

Hipotetik kıyasın geriye doğru ispatını Bölüm~2'de görmüştük:

\leandosya[Geriye doğru]{Bolum05}{kiyas_geri}

\noindent Aynı ispatı ileriye doğru da yazabiliriz. \leantac{have}
taktiği, bağlama yeni bir ara sonuç ekler:

\leandosya[İleriye doğru]{Bolum05}{kiyas_ileri}

\begin{hedefgecisi}
\begin{hedefdurumu}[\texttt{have q} öncesi]
h1 : P → Q
h2 : Q → R
p : P
⊢ R
\end{hedefdurumu}
\sonra
\begin{hedefdurumu}[\texttt{have q} sonrası]
h1 : P → Q
h2 : Q → R
p : P
q : Q
⊢ R
\end{hedefdurumu}
\end{hedefgecisi}

\noindent \leantac{have} hedefi değiştirmez; yalnızca bağlamı
büyütür. Matematik kitaplarındaki ``\ldots olduğundan \ldots elde
ederiz'' cümlelerinin karşılığı budur.

\paragraph{suffices: ``Şunu göstermek yeter''}
Geriye doğru akıl yürütürken ara hedefe bir ad vermek isteyebiliriz.
\leantac{suffices} tam bunu yapar: ``$Q$ elimde olsaydı işi şöyle
bitirirdim; o hâlde $Q$'yu göstermek yeter.''

\leandosya[Yeter koşulu]{Bolum05}{kiyas_suffices}

\paragraph{show: hedefi açıkça yazmak}
\leantac{show} taktiği hiçbir şeyi değiştirmez; hedefin ne olduğunu
okura (ve \Lean{}'e) açıkça söyler. Hedef yazdığınız şeyle uyuşmazsa
\Lean{} hata verir. Uzun ispatlarda \leantac{show}, ispatı okunabilir
kılmanın en ucuz yoludur.

\leandosya[Hedefi açıkça yazmak]{Bolum05}{kiyas_show}

\begin{dikkat}[Dikkat: Hangi yön ne zaman?]
Geriye doğru ispatlar kısadır ama okunması zordur: Okur, ara hedefleri
görmek için \Lean{}'i çalıştırmak zorunda kalır. İleriye doğru ispatlar
uzundur ama kâğıt üzerindeki matematiğe benzer. İyi bir ispat genellikle
ikisini karıştırır: Büyük adımları \leantac{have} ile adlandırır, küçük
boşlukları \leantac{apply} ve \leantac{exact} ile kapatır.
\end{dikkat}

% ---------------------------------------------------------------------
\section{Bir ispat stratejisi}\label{sec:strateji}

Önermeler mantığında bir ispata nereden başlayacağınızı çoğu zaman
\emph{hedefin} ve \emph{hipotezlerin} biçimi söyler. \Cref{fig:strateji}
bu bilgiyi bir harita olarak topluyor.

\begin{figure}[p]
\centering
\begin{tikzpicture}[
    bas/.style={draw=koyu, fill=koyu, text=white, rounded corners=4pt,
                font=\baslikfont\bfseries\small, minimum width=5.6cm,
                minimum height=0.8cm},
    sekil/.style={draw=#1, fill=#1!8, thick, rounded corners=3pt,
                  font=\leanfont\small, minimum width=1.7cm, minimum height=0.75cm},
    eylem/.style={draw=black!20, fill=white, rounded corners=3pt,
                  font=\small, text width=4.3cm, minimum height=0.75cm,
                  align=left, inner sep=4pt},
    ok/.style={-{Stealth[length=2.2mm]}, thick, black!45},
  ]
  % ----- Sol sütun: hedef -----
  \node[bas] at (3.1,0) {1. Hedefe bak: $\vdash$ ?};
  \foreach \i/\s/\e in {
      1/{⊢ A → B}/{\leantac{intro h}},
      2/{⊢ ¬A}/{\leantac{intro h}; hedef \lc{False} olur},
      3/{⊢ A ∧ B}/{\leantac{constructor}},
      4/{⊢ A ↔ B}/{\leantac{constructor}: iki yön},
      5/{⊢ A ∨ B}/{Taraf belliyse \lc{Or.inl}/\lc{Or.inr}; değilse önce 2. adım},
      6/{⊢ True}/{\leantac{trivial}},
      7/{⊢ False}/{Çelişki ara: \lc{hn h}, \leantac{contradiction}},
      8/{⊢ atom}/{2. adıma geç}}
  {
    \node[sekil=ikincil] (h\i) at (0.85,{-1.3*\i}) {\s};
    \node[eylem, anchor=west] (e\i) at (1.95,{-1.3*\i}) {\e};
    \draw[ok] (h\i) -- (e\i);
  }
  % ----- Sağ sütun: hipotezler -----
  \node[bas] at (10.1,0) {2. Hipotezlere bak: h : ?};
  \foreach \i/\s/\e in {
      1/{h : A ∧ B}/{\leantac{obtain ⟨a, b⟩ := h}},
      2/{h : A ∨ B}/{\leantac{rcases h with a | b}: iki durum},
      3/{h : A → B}/{Hedef \lc{B}: \leantac{apply h}; \lc{a : A} varsa: \leantac{have b := h a}},
      4/{h : ¬A}/{Hedef \lc{False}: \leantac{apply h}, sonra \lc{A}'yı göster},
      5/{h : A ↔ B}/{\leantac{rw [h]} ya da \lc{h.mp}, \lc{h.mpr}},
      6/{h : False}/{\leantac{exact h.elim}; bitti},
      7/{h : hedef}/{\leantac{exact h}, \leantac{assumption}}}
  {
    \node[sekil=vurgu] (k\i) at (7.85,{-1.3*\i}) {\s};
    \node[eylem, anchor=west] (f\i) at (8.95,{-1.3*\i}) {\e};
    \draw[ok] (k\i) -- (f\i);
  }
  % ----- Alt: takılınca -----
  \node[draw=teoremframe, fill=teorembg, thick, rounded corners=4pt,
        text width=13.6cm, align=left, font=\small, inner sep=6pt]
     at (6.6,-12.45) {%
     {\baslikfont\bfseries\color{teoremframe} 3. Takıldıysan:}
     Kâğıt üzerinde bir doğruluk tablosu ya da küçük bir Kripke modeli
     dene. Önerme klasik olarak doğru ama yapıcı olarak ispatlanamıyorsa
     (\cref{bol:yapici-klasik}), içinde bir ``hangisi?'' sorusu
     saklıdır: \lc{Classical.em} ile durum ayır. Önerme klasik olarak da
     yanlışsa, bir karşı örnek bul ve ispatlamaya çalışmayı bırak!};
\end{tikzpicture}
\caption{Önermeler mantığında ispat stratejisi. Önce hedefin biçimine,
sonra hipotezlerin biçimine bakılır. Her satır bir giriş (sol) ya da
eleme (sağ) kuralına karşılık gelir.}
\label{fig:strateji}
\end{figure}

Bu haritanın arkasında güçlü bir teorik sonuç yatar: Yapıcı önermeler
mantığında, hedefin biçimine göre giriş kurallarını, hipotezlerin
biçimine göre eleme kurallarını uygulayan ve aynı hedefe dönmemeye
dikkat eden sistemli bir arama, ispat varsa onu \emph{her zaman} bulur. Bu sonucun kesin hâli, \cref{bol:ardisik}'te
göreceğimiz kesme eleme teoreminin bir sonucudur. Yapıcı önermeler
mantığının \emph{karar verilebilir} olmasının nedeni de budur:
Herhangi bir önermenin ispatlanabilir olup olmadığını sonlu bir aramayla
anlayabiliriz.

% ---------------------------------------------------------------------
\section{Kapsamlı bir örnek}\label{sec:kapsamli-ornek}

Stratejiyi bir örnek üzerinde adım adım uygulayalım. İspatlamak
istediğimiz önerme şudur:
\[
  (P \to R) \wedge (Q \to R) \;\leftrightarrow\; (P \vee Q \to R).
\]
``$P$ de $R$'yi gerektiriyorsa $Q$ da gerektiriyorsa, $P$ veya $Q$
$R$'yi gerektirir; tersi de doğrudur.''

\begin{enumerate}[adimlar]
  \item \textbf{Hedef \lc{↔}}: \leantac{constructor} ile iki yöne ayır.
  \item \textbf{(→) yönü, hedef \lc{→}}: \leantac{intro} ile
        \lc{(P → R) ∧ (Q → R)}'yi al. Bu bir \lc{∧} hipotezi:
        \leantac{obtain} ile parçala. Hedef yine \lc{→}: \lc{P ∨ Q}'yu
        al. Bu bir \lc{∨} hipotezi: durum ayır. \leantac{rintro}, bu dört
        adımı tek satırda yapar.
  \item \textbf{Her durumda hedef \lc{R}} (atom): Hipotezlere bak.
        \lc{hpr : P → R} ve \lc{hp : P} var: \lc{hpr hp}.
  \item \textbf{(←) yönü}: \lc{h : P ∨ Q → R}'yi al. Hedef \lc{∧}: iki
        bileşeni ayrı ayrı kur. Her bileşen bir implikasyon: \lc{fun hp => \ldots}.
\end{enumerate}

\leandosya[Durum ayırmanın denkliği]{Bolum05}{durum_iff}

\noindent (→) yönünün doğal çıkarım ağacı \cref{fig:durum-agaci}'de
görülüyor. Ağaçtaki her kural, \Lean{} ispatındaki bir parçaya karşılık
geliyor.

\begin{figure}[htbp]
\centering
\small
\begin{prooftree}
  \hypo{[P \vee Q]^{4}}
  %
  \hypo{[(P \to R) \wedge (Q \to R)]^{3}}
  \infer1[$\wedge$E$_1$]{P \to R}
  \hypo{[P]^{1}}
  \infer2[$\to$E]{R}
  %
  \hypo{[(P \to R) \wedge (Q \to R)]^{3}}
  \infer1[$\wedge$E$_2$]{Q \to R}
  \hypo{[Q]^{2}}
  \infer2[$\to$E]{R}
  %
  \infer3[$\vee$E$^{1,2}$]{R}
  \infer1[$\to$G$^{4}$]{P \vee Q \to R}
  \infer1[$\to$G$^{3}$]{(P \to R) \wedge (Q \to R) \to (P \vee Q \to R)}
\end{prooftree}

\vspace{8pt}
\begin{tikzpicture}[font=\leanfont\scriptsize]
  \node[draw=black!20, rounded corners=3pt, fill=white, align=left, inner sep=6pt]{%
    \textcolor{leantac}{\bfseries rintro} ⟨hpr, hqr⟩ (hp | hq)
    \quad$\rightsquigarrow$\quad {\rmfamily\small $\to$G$^3$, $\wedge$E, $\to$G$^4$, $\vee$E$^{1,2}$}\\[2pt]
    \textcolor{leantac}{\bfseries exact} hpr hp \;/\; \textcolor{leantac}{\bfseries exact} hqr hq
    \quad$\rightsquigarrow$\quad {\rmfamily\small iki $\to$E adımı}};
\end{tikzpicture}
\caption{$(P \to R) \wedge (Q \to R) \to (P \vee Q \to R)$'nin doğal
çıkarım ağacı ve \Lean{} ispatıyla karşılaştırması.}
\label{fig:durum-agaci}
\end{figure}

\noindent Terim modunda aynı ispat iki satıra sığar. (→) yönünde
\lc{Or.elim h hpr hqr} ifadesi, $\vee$E kuralını \emph{doğrudan}
uygular: İki durumun her biri için hazır bir fonksiyonumuz zaten var.

\leandosya[Durum ayırmanın denkliği, terim modu]{Bolum05}{durum_iff_terim}

% ---------------------------------------------------------------------
\section{İspat normalleştirme}\label{sec:normallestirme}

Bazı ispatlar gereksiz dolambaçlar içerir. En basit örnek, bir ``ve''
kanıtını kurup hemen ardından sökmektir:

\leandosya[Dolambaçlı ve dolaysız ispat]{Bolum05}{dolambacli}

\noindent \lc{dolambacli} ispatı, \lc{⟨hp, hq⟩} çiftini kurar ($\wedge$G)
ve hemen sol bileşenini alır ($\wedge$E$_1$). Sonuçta elimize başta
zaten sahip olduğumuz \lc{hp} geçer. Ağaç olarak (\cref{fig:normallestirme}):

\begin{figure}[htbp]
\centering
\begin{tikzpicture}
  \node (sol) at (0,0) {%
    \begin{prooftree}
      \hypo{\mathcal{D}_1}\ellipsis{}{P}
      \hypo{\mathcal{D}_2}\ellipsis{}{Q}
      \infer2[\textcolor{ikincil}{$\wedge$G}]{P \wedge Q}
      \infer1[\textcolor{vurgu}{$\wedge$E$_1$}]{P}
    \end{prooftree}};
  \node (sag) at (5.2,0) {%
    \begin{prooftree}
      \hypo{\mathcal{D}_1}\ellipsis{}{P}
    \end{prooftree}};
  \draw[-{Stealth[length=3mm]}, very thick, anarenk] (1.9,0) -- node[above, font=\baslikfont\scriptsize\bfseries, text=anarenk]{indirgeme} (4.4,0);
  \node[font=\baslikfont\scriptsize, text=black!60] at (0,-1.6) {dolambaç: giriş, hemen ardından eleme};
  \node[font=\baslikfont\scriptsize, text=black!60] at (5.2,-1.6) {normal};

  \node (sol2) at (0,-4.1) {%
    \begin{prooftree}
      \hypo{[P]^{1}}\ellipsis{$\mathcal{D}$}{Q}
      \infer1[\textcolor{ikincil}{$\to$G$^1$}]{P \to Q}
      \hypo{\mathcal{E}}\ellipsis{}{P}
      \infer2[\textcolor{vurgu}{$\to$E}]{Q}
    \end{prooftree}};
  \node (sag2) at (5.2,-4.1) {%
    \begin{prooftree}
      \hypo{\mathcal{E}}\ellipsis{}{P}
      \ellipsis{$\mathcal{D}$}{Q}
    \end{prooftree}};
  \draw[-{Stealth[length=3mm]}, very thick, anarenk] (1.9,-4.1) -- node[above, font=\baslikfont\scriptsize\bfseries, text=anarenk]{indirgeme} (4.4,-4.1);
  \node[font=\baslikfont\scriptsize, text=black!60, align=center] at (5.2,-6.0) {$[P]$ varsayımının yerine\\$\mathcal{E}$ türetimi konur};
\end{tikzpicture}
\caption{İki temel indirgeme. Bir giriş kuralının hemen ardından aynı
bağlacın eleme kuralı geliyorsa, ikisi birlikte silinebilir.}
\label{fig:normallestirme}
\end{figure}

\begin{tanim}{Normal türetim}{normal}
Bir türetimde, bir bağlacın giriş kuralıyla elde edilen bir önermeye
hemen ardından aynı bağlacın eleme kuralı uygulanıyorsa, bu önermeye
\textbf{dolambaç} (İng.\ \emph{detour}, \emph{maximal formula}) denir.
Dolambaç içermeyen türetimlere \textbf{normal} türetim denir.
\end{tanim}

İmplikasyondaki dolambaç da aynı fikri izler: Önce bir fonksiyon kurar,
sonra onu hemen bir argümana uygularız.

\leandosya[İmplikasyonda dolambaç]{Bolum05}{dolambacli_imp}

\noindent Bu ispatın indirgenmiş hâli \lc{hpq hp}'dir: Fonksiyonun
gövdesinde \lc{h} gördüğümüz her yere \lc{hp} koyarız. Programlama
dillerinde buna \textbf{$\beta$-indirgeme}, yani \emph{bir fonksiyonu
çağırmak} denir. Aynı dolambaçları veri dünyasında yazdığımızda
\Lean{} onları \emph{hesaplar}:

\leandosya[Normalleştirme = hesaplama]{Bolum05}{hesaplama}

\begin{teorem}{Normalleştirme, Prawitz 1965}{normallestirme}
Yapıcı önermeler mantığında her türetim, aynı varsayımlardan aynı sonucu
veren bir normal türetime indirgenebilir \citep{prawitz1965}.
\end{teorem}

\begin{onemlifikir}[Önemli Fikir: Normalleştirme = program çalıştırma]
Curry--Howard eşlemesi altında Prawitz'in teoremi şunu söyler: Her
ispat (program) dolambaçları silinerek (hesaplanarak) bir sonuca ulaşır;
hesaplama sonsuza dek sürmez. İspat kuramında ``normalleştirme'', bilgisayar
biliminde ``program çalıştırma'' denen şeyin aynısıdır. Bu bağlantıyı
\cref{bol:lambda}'te $\lambda$-hesabı ile kesinleştireceğiz.
\end{onemlifikir}

Normal türetimlerin önemli bir özelliği vardır: \emph{Alt formül özelliği}.
Normal bir türetimde görünen her önerme, ya sonucun ya da açık
varsayımlardan birinin bir parçasıdır. Yani normal bir ispat ``dışarıdan''
hiçbir yeni kavram getirmez. Bundan çarpıcı bir sonuç çıkar:

\begin{teorem}{Tutarlılık}{tutarlilik}
Yapıcı önermeler mantığında $\bot$ önermesi, varsayımsız olarak
türetilemez: $\nvdash \bot$.
\end{teorem}

\noindent\textit{İspat taslağı.} $\bot$'un bir türetimi olsaydı,
Teorem~\ref{trm:normallestirme}'e göre normal bir türetimi de olurdu. Alt
formül özelliğine göre bu türetimdeki her önerme $\bot$'un bir parçası,
yani $\bot$'un kendisi olmalıdır. Ama $\bot$'un giriş kuralı yoktur ve
açık varsayım da yoktur; dolayısıyla ağacın hiçbir yaprağı olamaz.
Bu bir çelişkidir. \hfill\qedsymbol

\noindent Bu, mantığın \emph{kendi tutarlılığını} kendi araçlarıyla
gösteren ilk ve en basit örnektir. Aynı fikri aritmetiğe uygulama
çabası, Gentzen'i 1936'da aritmetiğin tutarlılık ispatına götürdü
\citep{gentzen1936}; bu hikâyeye \ciltiki{4}'te döneceğiz.

% ---------------------------------------------------------------------
\section{Taktik birleştiricileri}\label{sec:birlestiriciler}

Bazı ispatlarda aynı taktiği birçok hedefe uygulamamız gerekir. Üç
önermenin ``ve''sini kurmak için:

\leandosya[Uzun yol]{Bolum05}{birlestirici_uzun}

\noindent Bu ispatı \emph{taktik birleştiricileri} ile kısaltabiliriz:

\leandosya[Kısa yol]{Bolum05}{birlestirici_kisa}

\begin{center}
\renewcommand{\arraystretch}{1.45}
\rowcolors{2}{tablosatir}{white}
\begin{tabularx}{\linewidth}{>{\raggedright\arraybackslash}p{3.6cm} X}
\rowcolor{tabloust}
\tablobaslik{Birleştirici} & \tablobaslik{Ne yapar?} \\
\lc{t₁ <;> t₂}                     & \lc{t₁}'i uygular, \emph{ortaya çıkan her hedefe} \lc{t₂}'yi uygular. \\
\leantac{all\_goals} \lc{t}        & Açık olan bütün hedeflere \lc{t}'yi uygular. \\
\leantac{first} \lc{| t₁ | t₂}     & Önce \lc{t₁}'i dener; başarısız olursa \lc{t₂}'yi. \\
\leantac{try} \lc{t}               & \lc{t}'yi dener; başarısız olursa hiçbir şey yapmaz. \\
\leantac{repeat} \lc{t}            & \lc{t}'yi başarısız olana dek tekrar tekrar uygular. \\
\leantac{refine} \lc{⟨?\_, ?\_⟩}   & Terimin bir kısmını yazar; \lc{?\_} ile işaretli boşlukları yeni hedef yapar. \\
\end{tabularx}
\end{center}

\leantac{refine}, terim modu ile taktik modunu birleştirir: İspatın
iskeletini bir terim olarak yazar, boşlukları taktiklerle doldururuz.
\leantac{first} ile durum ayırmayı birleştiren bir örnek:

\leandosya[Durum ayırma ve first]{Bolum05}{birlestirici_veya}

\noindent İki durumun her birinde \Lean{} önce \lc{Or.inr h}'yi dener;
türler uyuşmazsa \lc{Or.inl h}'ye geçer.

\begin{dikkat}[Dikkat: Kısa her zaman iyi değildir]
Birleştiriciler ispatı kısaltır ama okunmasını zorlaştırabilir. Bir
ispat \lc{<;>} ile bir satıra indiğinde, okur hangi hedefin hangi
taktikle kapandığını göremez. Bir kural olarak, birleştiricileri
\emph{aynı biçimdeki} hedefler için kullanın (\lc{refine ⟨?\_, ?\_, ?\_⟩
<;> assumption} gibi), farklı mantık gerektiren durumlar için değil.
\end{dikkat}

\paragraph{Otomasyona bir bakış}
Önermeler mantığı karar verilebilir olduğuna göre, bu bölümdeki
ispatları otomatik olarak bulan bir taktik yazılabilir. \Mathlib{}'teki
\leantac{tauto} taktiği tam olarak bunu yapar. \Lean{}'in güncel
sürümlerinde \lc{exact?} komutu ise kütüphanede hedefi kapatan bir
teorem arar. Bu araçları \ciltiki{6}'da tanıyacak,
\ciltiki{7}'de ise kendi önermeler mantığı ispatçımızı
yazacağız. O zamana dek ispatları elle yazmanızı öneririz: Otomasyon,
ne yaptığını anlayan birinin elinde güçlü bir araçtır.

% ---------------------------------------------------------------------
\begin{tarihce}[Tarihçe: Hilbert sistemlerinden Gentzen'e]
Frege'den Hilbert'e kadar mantık sistemleri \emph{aksiyomatik} olarak
kurulurdu: Birkaç aksiyom şeması (örneğin
$P \to (Q \to P)$ ve $(P \to Q \to R) \to (P \to Q) \to P \to R$) ve
tek bir kural (modus ponens). Bu \emph{Hilbert sistemleri} teorik olarak
güçlüdür, ama içlerinde ispat yazmak zahmetlidir: $P \to P$ gibi basit
bir önermenin ispatı bile beş adım tutar. Gentzen 1935'te bu sistemlerin
``gerçek matematiksel akıl yürütmeden çok uzak'' olduğunu yazarak doğal
çıkarımı önerdi \citep{gentzen1935}. Bugün ispat asistanlarının
neredeyse tamamı doğal çıkarım ya da onun bir uzantısı üzerine kuruludur.
Hilbert sisteminin aksiyomlarının Bölüm~2'deki $K$ ve $S$
kombinatörleri olduğuna dikkat edin.
\end{tarihce}

\begin{arastirmanotu}
\textbf{Uyum ve ispat-kuramsal anlambilim.} \Cref{sec:dogal-cikarim-fikri}'te
sözünü ettiğimiz giriş ve eleme kuralları arasındaki \emph{uyum}, Michael
Dummett ve Dag Prawitz'in \emph{ispat-kuramsal anlambilim} programının
temel kavramıdır: Bir bağlacın anlamı, doğruluk tablosuyla değil, giriş
kurallarıyla verilir; eleme kuralları bu anlamla uyumlu olmalıdır.
Normalleştirme teoremi, bu uyumun teknik ifadesidir
\citep{prawitz1965, dummett2000}.

\textbf{Taktik dilleri.} \Lean{}'in taktik dili, kendisi \Lean{}'de
yazılmış bir programlama dilidir; kullanıcılar yeni taktikler
tanımlayabilir \citep{demoura2021lean4, ullrich2022macros}. Taktik
ispatlarının okunabilirliği ve bakımı, büyük formalizasyon projelerinde
pratik bir sorun olarak tartışılmaya devam ediyor; \Mathlib{} bu nedenle
belirli bir ispat stili önerir (\ciltiki{9}).

\textbf{İleri okuma.} Normalleştirme ve alt formül özelliğinin tam
ispatları için \citet{troelstra2000}; Curry--Howard açısından
normalleştirme için \citet{girard1989} ve \citet{sorensen2006}.
\end{arastirmanotu}

\begin{bolumozeti}
\begin{itemize}
  \item Her bağlacın giriş ve eleme kuralları vardır; tek klasik kural
        RAA'dır (\cref{fig:kural-katalogu}).
  \item \Lean{}'in hedef ekranı bir ardışıktır: $\Gamma \vdash \varphi$.
        Bir taktik, bir kuralı aşağıdan yukarı uygular.
  \item İspatlar geriye (\leantac{apply}, \leantac{suffices}) ya da
        ileriye (\leantac{have}) doğru kurulabilir; iyi ispatlar ikisini
        karıştırır.
  \item Strateji: Önce hedefin biçimine (giriş kuralları), sonra
        hipotezlerin biçimine (eleme kuralları) bak.
  \item Dolambaç, bir girişin hemen ardından gelen elemedir. Normalleştirme
        dolambaçları siler; Curry--Howard altında bu, program çalıştırmaktır.
        Normal ispatların alt formül özelliği, tutarlılığı verir.
  \item \lc{<;>}, \leantac{all\_goals}, \leantac{first} ve \leantac{refine}
        ispatları kısaltır; ölçülü kullanılmalıdır.
\end{itemize}
\end{bolumozeti}

\begin{alistirmalar}
\alistirma{1} Varsayımların sırasını değiştirin:
\lc{(P → Q → R) → (Q → P → R)}. Terim modunda bu ispat hangi
fonksiyondur? (\lc{a5\_1})

\alistirma{1} \lc{(P ∧ Q → R) → P → Q → R} önermesini önce yalnızca
\leantac{have} ve \leantac{exact} ile (\lc{a5\_2}), sonra yalnızca
\leantac{apply} ve \leantac{exact} ile (\lc{a5\_3}) ispatlayın. Hangisi
size daha okunabilir geldi?

\alistirma{2} İmplikasyonun ve'ye dağıldığını ispatlayın:
\lc{(P → Q ∧ R) ↔ (P → Q) ∧ (P → R)}. (\lc{a5\_4})

\alistirma{2} \lc{(P ∨ Q) ∧ (P → R) ∧ (Q → R) → R}. Önce
\cref{fig:strateji}'yi kullanarak bir plan yazın, sonra ispatlayın.
(\lc{a5\_5})

\alistirma{2} \lc{P ∧ Q ∧ R → R ∧ Q ∧ P} önermesini \leantac{rintro}
ve tek bir \lc{refine ... <;> assumption} satırıyla ispatlayın.
(\lc{a5\_6})

\alistirma{2} Dört halkalı zinciri önce geriye, sonra ileriye doğru
ispatlayın: \lc{(P → Q) → (Q → R) → (R → S) → P → S}. (\lc{a5\_7})

\alistirma{3} \lc{(P ∨ Q → R) → ¬R → ¬P ∧ ¬Q}. Doğal çıkarım ağacını
da çizin. (\lc{a5\_8})

\alistirma{3} \lc{(P ∨ Q) ∧ (R ∨ S) → (P ∧ R) ∨ (P ∧ S) ∨ (Q ∧ R) ∨ (Q ∧ S)}.
İpucu: \leantac{rintro} ile \lc{⟨hp | hq, hr | hs⟩} örüntüsü dört durumu
birden açar. (\lc{a5\_9})

\alistirma{2} (Kâğıt-kalem) Aşağıdaki türetimi normalleştirin. Kaç
indirgeme adımı gerekti?
\[
\begin{prooftree}
  \hypo{[P]^{1}}
  \hypo{[Q]^{2}}
  \infer2[$\wedge$G]{P \wedge Q}
  \infer1[$\wedge$E$_2$]{Q}
  \infer1[$\to$G$^{1}$]{P \to Q}
  \hypo{P}
  \infer2[$\to$E]{Q}
  \infer1[$\to$G$^{2}$]{Q \to Q}
\end{prooftree}
\]

\alistirma{3} (Kâğıt-kalem) Alt formül özelliğini kullanarak $P \vee Q$'nun
varsayımsız olarak türetilemeyeceğini gösterin ($P$, $Q$ farklı
atomlar). İpucu: Normal bir türetimin son kuralı ne olabilir?

\alistirma{3} (Tartışma) \cref{fig:strateji}'deki haritayı kullanarak
Peirce yasasını ispatlamaya çalışın. Haritanın hangi adımında
takılıyorsunuz? Bu takılma, \cref{bol:yapici-klasik}'teki Kripke
modeliyle nasıl ilişkilidir?
\end{alistirmalar}
